% Figure: hot and cold composite curves of pinch analysis, shifted along the
% enthalpy axis until the minimum vertical gap equals the chosen minimum approach;
% the overlap is the heat that can be recovered, the overshoots are the minimum
% hot and cold utility, and the closest approach is the pinch.
\begin{figure}[htbp]
\centering
\begin{tikzpicture}
\begin{axis}[
  width=12cm, height=7.5cm,
  xlabel={cumulative enthalpy $H$ (MW)},
  ylabel={temperature ($^{\circ}$C)},
  xmin=0, xmax=10, ymin=20, ymax=200,
  axis lines=left,
  legend style={font=\scriptsize, at={(0.02,0.98)}, anchor=north west},
  legend cell align=left,
  tick label style={font=\scriptsize},
  label style={font=\small},
  clip=false,
]
% hot composite curve (heat available from hot streams), shifted to the right
\addplot[engred, very thick] coordinates {(2.0,55) (4.0,90) (7.0,130) (9.0,180)};
\addlegendentry{hot composite}
% cold composite curve (heat required by cold streams), to the left
\addplot[engblue, very thick] coordinates {(0.0,30) (2.5,75) (5.5,120) (7.5,160)};
\addlegendentry{cold composite}
% pinch: closest vertical approach near H = 4.5, mark with a thin vertical bar
\draw[enggray, thick, <->] (axis cs:4.2,98) -- (axis cs:4.2,108);
\node[enggray, font=\scriptsize, anchor=west] at (axis cs:4.3,103) {$\Delta T_{\min}$ (pinch)};
% recoverable-heat overlap bracket along the enthalpy axis
\draw[engamber, thick, |-|] (axis cs:2.0,33) -- (axis cs:7.5,33);
\node[engamber, font=\scriptsize, anchor=north] at (axis cs:4.75,33) {heat recovered};
% minimum cold utility (cold-end overshoot of hot curve)
\draw[engblue, thick, |-|] (axis cs:0.0,25) -- (axis cs:2.0,25);
\node[engblue, font=\scriptsize, anchor=north] at (axis cs:1.0,25) {min cold utility};
% minimum hot utility (hot-end overshoot of cold curve)
\draw[engred, thick, |-|] (axis cs:7.5,190) -- (axis cs:9.0,190);
\node[engred, font=\scriptsize, anchor=south] at (axis cs:8.25,190) {min hot utility};
\end{axis}
\end{tikzpicture}
\caption[Composite curves and the pinch]{Hot and cold composite curves of a
process, each built by adding the enthalpy of all streams within a temperature
interval, then slid along the enthalpy axis until the closest vertical gap equals
the chosen minimum approach $\Delta T_{\min}$. The horizontal overlap is the heat
that internal exchangers can recover; the cold-end overshoot of the hot curve is
the minimum cold utility the process still needs, and the hot-end overshoot of the
cold curve is the minimum hot utility. The point of closest approach is the pinch,
the bottleneck that fixes how much heat integration the process allows. Designing
across the pinch, by transferring heat from above it to below it, wastes utility
on both sides and is the error the method exists to prevent.}
\label{fig:vol04ch06:composite-curves}
\end{figure}
